\section{Compatibilidad de caudal, recarga y retirada por gas (GAS-28)}
\label{sec:sd4-integracion-gas28}
La renovación instalada de 900 m$^3$/h efectivos por banco se distingue del permiso de recarga. Con emisión conservadora de electrólisis de 7,5978 m$^3$/h a 60 A agregados, ya desarrollada para 240 celdas en serie, la mezcla ideal estacionaria resulta $x=q/Q=0,008442$, 0,8442\,\% H$_2$. Supera el umbral propio de intervención de 0,8\,\% y la prealarma de 0,4\,\%. Por ello no se autoriza recarga permanente a 60 A con ese caudal: el corte eventual no convierte el régimen en compatible. Los 900 m$^3$/h siguen como punto de selección ambiental y de purga de la arquitectura, sujeto a medición efectiva, sin describirse como renovación suficiente para mantener libre de alarma aquel extremo de corriente \parencites{p7-eaton-93t-guia}{csbManualJunio2025}{nistCODATA2022}.

Para no superar 0,4\,\% con ese mismo modelo se requiere $Q\geq7,5978/0,004=1899,45$ m$^3$/h por banco, antes de incertidumbre y reparto; para 0,8\,\%, $Q\geq949,725$ m$^3$/h. Ninguna de esas desigualdades acredita un caudal real seguro ni una selección de extracción, reposición, serpentín y energía a 1.900 m$^3$/h. Elevar sólo un ventilador invalidaría humedad, ductos, velocidades, ruido, fuentes y autonomía. La oferta conserva la arquitectura de 900 m$^3$/h y exige un límite efectivo de recarga, autorizado y supervisado por cadena, antes de habilitarla. No se prescribe subir umbrales para acomodar el cálculo.

La proporcionalidad conservadora con corriente permite calcular, exclusivamente para cribado ideal, $I\leq60(0,004)(900)/7,5978=28,4293$ A agregados. Ese valor es un límite sin margen, no un ajuste de cargador seleccionado. El ajuste requerido de 20 A daría 0,2814\,\% estacionario ideal; necesita aún perfil OEM autorizado, control térmico/tensión, desigualdad entre cadenas, margen de medida, caudal distribuido e interlock efectivo. No se traslada automáticamente a un cargador sustituto ni se da por satisfecho por una opción de pantalla configurable. El permiso permanece retirado mientras esa cadena de condiciones no esté individualizada. La retirada de rectificador por MN elimina recarga AC pero lleva el camino a batería; exige gestionar continuidad y reserva, sin ciclar MN como regulador \parencites{p7-eaton-93t-guia}{csbHRL12540WRA260131}.

\textbf{Dinámica propia para dimensionar respuesta.} Bajo volumen de 36 m$^3$, caudal 900 m$^3$/h y mezcla perfecta constante, $dx/dt=q/V-(Q/V)x$ y $x(t)=q/Q+(x_0-q/Q)\exp(-Qt/V)$, con $t$ en horas. La constante de tiempo es 2,4 min. Desde aire ideal inicialmente sin gas y manteniendo la emisión de 60 A, alcanza 0,4\,\% en aproximadamente 1,541 min; de 0,4 a 0,8\,\%, unos 5,538 min. Sin extracción, entre esos dos umbrales sólo dispone de aproximadamente 68,23 s. Son tiempos de modelo, no latencias permitidas de sensor ni ensayos de dispersión: estratificación, volumen desplazado por equipos y bolsas ciegas pueden acortar la ventana local. Se exige retirar permiso ante pérdida de caudal/señal válida sin esperar a que H$_2$ cruce alarma. El presupuesto de latencia incluye detector/filtro, transporte, lógica, contacto, caída de bobina y apertura del elemento final; no se confunde con los 39/139 ms parciales de CEN-26.

\textbf{Límite del sellado por incendio.} Si el recinto comienza a 0,4\,\%, mantenerlo por debajo de 1,0\,\% durante diez minutos sin extracción impondría $q_{res}\leq(0,010-0,004)(36)/(10/60)=1,296$ m$^3$/h con el mismo modelo. Es apenas 17,06\,\% de la emisión conservadora precedente. La medición baja previa al cierre y el corte AC no demuestran ese límite residual. Tampoco el agente que permanece modifica por sí solo la generación de gas ni garantiza oxígeno suficiente para una lectura catalítica válida. GAS-28 fija esta interfaz cuantificada para la selección gas/incendio: retención de agente, emisiones posaislamiento y maniobra aprobada deben resolverse conjuntamente antes de habilitar sellado o descarga en banco. Los sectores ordinarios conservan sus secuencias independientes.
